\documentclass[10pt,a4paper]{article}
\usepackage[margin=25mm]{geometry}
\usepackage{amsmath,amssymb,amsthm,hyperref}
\setlength{\emergencystretch}{2em}
\newtheorem{proposition}{Proposition}
\newcommand{\gen}[1]{\langle#1\rangle}
\newcommand{\Alt}{\operatorname{Alt}}
\title{The complete supported-interval rank step}
\author{Human Compression project: bounded L044 proof component}
\date{}
\begin{document}\maketitle
Let $G=\gen{\rho_0,\ldots,\rho_{r-1}}\leq A_n$, $n\geq12$,
be transitive and a string C-group, with all maximal omissions
intransitive. Suppose the unique eligible interior $i$-edge
joins its two omission orbits $A,B$, and the active label sets
on these two orbits are each single intervals.
As proved in the accepted supported-interval lemma, after
interchanging the orbits the actual groups
$L=G_{<i}$ and $R=G_{>i}$ are supported on and transitive
on $A$ and $B$, respectively. Their even string C-markings
have connected diagrams, intransitive omissions and endpoint
anchors. Assume the alternating rank bound at every strictly
smaller degree at least $12$.

\begin{proposition}
Unless the actual first three or last three generators are one
of the two degree-six $A_5$ markings of source Table~2,
rows~(2)--(3), up to full string duality, one has
\[
                     r\leq\left\lfloor\frac{n-1}{2}\right\rfloor.
\]
\end{proposition}
\begin{proof}
Suppose the asserted rank bound fails. Write
$m_A=|A|,m_B=|B|$ and let $d_A,d_B$ be the two interval
ranks, so $r=d_A+d_B+1$. Both degrees are at least four,
since a nontrivial even involution cannot be supported on
fewer points.
For a side of degree $m<12$, the accepted complete small
anchored-imprimitive bound gives $d\leq\lfloor(m-2)/2\rfloor$
in the imprimitive case. In the primitive case, the accepted
proper-interval exclusion leaves the same bound unless the
marking is source Table~2 row~(1), (2) or (3).
Rows~(2)--(3) are the excluded actual three-generator
prefixes or suffixes. Row~(1) is the degree-five dihedral
head: the complete dihedral attachment proposition,
with its explicit six-choice normalization, gives the
required global rank bound or exactly those excluded
prefixes. Thus this case is impossible under our current
contrary assumptions.

For a side of degree $m\geq12$ which is proper in its
natural alternating group, the accepted primitive bound
and the complete anchored-imprimitive bound again give
$d\leq\lfloor(m-2)/2\rfloor$. For a full natural alternating
side the explicit smaller-degree induction hypothesis gives
$d\leq\lfloor(m-1)/2\rfloor$.
Consequently, unless both sides are full natural alternating
groups of degrees at least $12$,
\[
 r\leq\frac{m_A-2}{2}+\frac{m_B-1}{2}+1
     =\frac{n-1}{2},
\]
a contradiction. (Interchange the sides to put the strict
bound first.)

In the remaining case use the accepted faithful even
one-point extension lemma. Its full proof writes the
crossing generator as $(a\ b)uv$, chooses the even exterior
factor, and uses actual ambient IP to show that the
one-point projection kernel is trivial. Thus one of the
two sides, say $A$, extends to a genuine even string C-group
of degree $m_A+1$ and rank $d_A+1$, with an endpoint anchor
and all omissions intransitive. The complete primitive,
anchored-imprimitive and smaller-degree alternating bounds
apply to this \emph{faithful} marking and give
\[
 d_A+1\leq\left\lfloor\frac{m_A}{2}\right\rfloor,
 \qquad d_B\leq\left\lfloor\frac{m_B-1}{2}\right\rfloor.
\]
Since $m_A+1<n$, this induction is acyclic. Adding yields
the required rank bound and the final contradiction.
\end{proof}

\paragraph{Exact source and proof scope.}
This consolidates Propositions~5.15--5.17 of
Cameron--Fernandes--Leemans--Mixer, \emph{Highest rank of a
polytope for $A_n$}, arXiv:1605.09173v1, physical page~29.
The source's interval-deficit and one-point-extension strategy
is reused. All small degrees are now covered using the
accepted complete Table~2 and anchored bounds and the new
complete dihedral adapter. The full one-point kernel proof
remains in the paid common component; no projected IP is
assumed. The two exceptional $A_5$ prefixes and all later
mixed-support and connected-square source branches remain
separate. The induction hypothesis is explicit, so this note
alone does not prove the complete L044 theorem or admit an
original. No new computation or formal theorem is used.
\end{document}
